\documentclass[11pt,a4paper]{article}

\usepackage[T1]{fontenc}
\usepackage[utf8]{inputenc}
\usepackage[brazil]{babel}
\usepackage{lmodern}
\usepackage{microtype}
\usepackage[margin=2cm]{geometry}
\usepackage{hyperref}
\usepackage{titlesec}
\usepackage{xcolor}

\hypersetup{
  colorlinks=true,
  linkcolor=black,
  urlcolor=black,
  citecolor=black,
  pdfborder={0 0 0}
}

\definecolor{sectioncolor}{HTML}{111111}
\definecolor{accentcolor}{HTML}{1A1A1A}

\pagestyle{empty}
\setlength{\parindent}{0pt}
\setlength{\parskip}{5pt}
\linespread{1.05}
\selectfont

\titleformat{\section}
  {\bfseries\large\color{sectioncolor}}
  {}{0pt}{}[\vspace{-2pt}{\color{sectioncolor}\titlerule[0.8pt]}]
\titlespacing*{\section}{0pt}{12pt}{8pt}

\newcommand{\cvrole}[4]{%
  \vspace{4pt}%
  {\bfseries\color{accentcolor}#1}\hfill{\itshape #2}\\[-2pt]%
  {\itshape #3}\hfill #4\\[6pt]%
}

\begin{document}

\begin{center}
  {\fontsize{20}{24}\selectfont\bfseries Kamille Hartmann Maccagnan}\\[8pt]
  {\large Analista de Projetos \textbar{} PMO \textbar{} Gestão de Cronogramas}\\[10pt]
  Campo Comprido, Curitiba-PR \enspace\textbullet\enspace (41) 9 9928-4424 \enspace\textbullet\enspace
  \href{mailto:kamillehartmannm@gmail.com}{kamillehartmannm@gmail.com}\\[3pt]
  \href{https://linkedin.com/in/kamille-hartmann-maccagnan/}{linkedin.com/in/kamille-hartmann-maccagnan/}
\end{center}

\vspace{6pt}

\section{Resumo Profissional}

Profissional com experiência em apoio a projetos, PMO e acompanhamento de cronogramas, prazos e entregas em ambientes corporativos e multinacionais. Atua na organização de informações, elaboração de relatórios executivos, documentação de processos e participação em reuniões de status com registro de encaminhamentos. Vivência em controle de portfólio de projetos, análise de riscos, padronização de documentação e interface com stakeholders de diferentes áreas. Conhecimento em metodologias ágeis e Boas Práticas PMOBOK, com domínio de Excel, PowerPoint e ferramentas de gestão de demandas. Perfil organizado, analítico e colaborativo, com atenção a prazos, escopo e qualidade das entregas.

\section{Experiência Profissional}

\cvrole{Anjuss}{Analista de Suporte de TI}{07/2025 -- Atual}

Atuo na organização de demandas, padronização de processos e documentação operacional, com criação e atualização de POPs para a rede de lojas e equipe de T.I., garantindo rastreabilidade e consistência das informações. Lidero treinamentos de novos funcionários e apoio a atividades administrativas e analíticas ligadas à operação. Desenvolvo relatórios e dashboards em Power BI para acompanhamento de indicadores e desempenho de processos, apoiando a identificação de pendências e oportunidades de melhoria. Participo da comunicação entre usuários e áreas técnicas, registrando ocorrências e encaminhamentos para resolução.

\cvrole{HORSE do Brasil}{Estagiária de TI (IS/IT LATAM)}{07/2024 -- 07/2025}

Atuei no apoio a projetos e PMO em ambiente multinacional, com acompanhamento de portfólio de projetos, cronogramas, prazos, entregas e responsáveis por iniciativas estratégicas. Atualizei o andamento de projetos por meio de relatórios executivos e dashboards em Power BI, apoiando reuniões de status e a gestão de lideranças. Participei do levantamento de riscos, documentação de processos, análise de pendências e melhoria contínua, seguindo práticas de governança de projetos. Utilizei o ServiceNow para gestão de demandas e atuei como ponte entre áreas de negócio, fornecedores e stakeholders internacionais, participando de reuniões em até três idiomas e registrando alinhamentos e encaminhamentos.

\cvrole{Restaurante Familiar}{Analista de Dados e Suporte Técnico}{01/2023 -- 07/2024}

Atuei em atividades analíticas e administrativas para apoio à gestão do negócio, com elaboração de relatórios, indicadores de desempenho e dashboards em Power BI para acompanhamento de resultados operacionais e financeiros. Organizei e padronizei informações para maior confiabilidade dos dados e apoio ao controle de performance. Atuei no alinhamento entre áreas operacionais e administrativas, contribuindo para o cumprimento de rotinas, prazos e melhoria da eficiência dos processos.

\section{Competências Técnicas}

\textbf{Gestão de Projetos e PMO:} Apoio a planejamento de projetos, acompanhamento de cronogramas e entregas, controle de prazos e escopo, portfólio de projetos, análise de riscos, reuniões de status, documentação de projetos, Boas Práticas PMOBOK, Scrum, Kanban.

\textbf{Relatórios e Comunicação:} Relatórios executivos e gerenciais, apresentações, Excel, PowerPoint, Power BI, comunicação com stakeholders, registro de decisões e encaminhamentos, trabalho em equipe.

\textbf{Organização e Processos:} Controle de informações, padronização de processos, POPs, gestão de demandas, melhoria contínua, atenção a detalhes e qualidade das entregas.

\textbf{Ferramentas:} ServiceNow, Trello, Monday, Power BI, Google Sheets, Pacote Office.

\textbf{Idiomas:} Inglês Avançado \enspace\textbullet\enspace Espanhol Básico \enspace\textbullet\enspace Coreano Básico

\section{Projetos e Conquistas}

1º lugar no Transformation Day Renault 2023 com solução orientada a dados e melhoria de processos. Participação no Hackathon Bosch 2023 com foco em resolução de problemas de negócio. Acompanhamento de portfólio de projetos e elaboração de relatórios executivos em ambiente multinacional na HORSE. Estruturação de KPIs e dashboards para monitoramento de performance. Desenvolvimento de documentações e POPs com padronização e rastreabilidade.

\section{Formação Acadêmica}

\textbf{PUCPR} \hfill Curitiba, PR\\
Graduação em Gestão da Tecnologia da Informação \hfill Previsão de conclusão: 06/2026\\[6pt]

\textbf{Escola Conquer} \hfill Curitiba, PR\\
Pós-graduação em Liderança e Gestão em Tecnologia \hfill Conclusão: 2024\\[6pt]

\textbf{Universidade Tuiuti do Paraná} \hfill Curitiba, PR\\
Graduação em Medicina Veterinária \hfill 06/2019\\[6pt]

\textbf{Qualittas} \hfill Curitiba, PR\\
Pós-graduação em Clínica Médica e Cirúrgica de Animais Exóticos e Pets Não Convencionais \hfill 12/2022

\end{document}
